% Part: set-theory
% Chapter: z
% Section: milestone

\documentclass[../../../include/open-logic-section]{subfiles}

\begin{document}

\olfileid{sth}{z}{milestone}
\olsection{$\Z^-$: একটি মাইলফলক}

পরের পরিচ্ছেদে \stagesinf{} নীতিতে আবার ফিরব। কিন্তু
অসীমতার স্বতঃসিদ্ধের সঙ্গে আমরা ইতিমধ্যে একটি গুরুত্বপূর্ণ
মাইলফলকে পৌঁছেছি। $\Zminus$ তত্ত্বের জন্য প্রয়োজনীয় সব
স্বতঃসিদ্ধ এখন আমাদের হাতে আছে। বিস্তারিতভাবে:

\begin{defn}
$\Zminus$ তত্ত্বের স্বতঃসিদ্ধগুলি হলো: বিস্তারগততা,
সংযোগ, জোড়, ঘাত সেট ও অসীমতার স্বতঃসিদ্ধ এবং
পৃথকীকরণ প্রকল্পের প্রতিটি নিদর্শন।
\end{defn}

নামটি \emph{জার্মেলো} সেটতত্ত্ব বোঝায় (পরে আলোচ্য
কিছু একটি \emph{বাদ দিয়ে})। জার্মেলোর নামেই এর পরিচয়
যথার্থ, কারণ তিনি তাঁর \citeyear{Zermelo1908Untersuchungen}-এর
কাজে মূলত এই তত্ত্বটি প্রণয়ন করেছিলেন।\footnote{ইতিহাস
ও কারিগরি দিক নিয়ে আকর্ষণীয় মন্তব্যের জন্য
\citet[Appendix
A]{Potter2004} দেখো।}

এই তত্ত্বে বিপুল পরিমাণ গণিত করা যায়। বিশেষ করে,
\olref[sfr][][]{part}-এ আগে যে কাজগুলি স্বতঃস্ফূর্তভাবে
করেছিলাম, সেগুলি যে $\Zminus$-এর মধ্যে আরও বিধিবদ্ধভাবে
করা যায়, ফিরে দেখে তা \emph{নিজেকে বোঝানো উচিত}।
(কিছুটা যাচাই করার পর চাইলে এগিয়ে গিয়ে
\olref[sth][z][arbintersections]{sec} পড়তে পারো।)
তাই এখন থেকে, আর আলাদা করে না বললেও, আমরা ধরে নেব
যে আমরা (অন্তত) $\Zminus$-এ কাজ করছি।

\end{document}
